\documentclass[11pt]{article}
\usepackage{lmodern}
\usepackage[T1]{fontenc}
\usepackage{amsmath,amssymb,booktabs,longtable,array,calc}
\usepackage[margin=1in]{geometry}
\usepackage[unicode,hidelinks]{hyperref}
\providecommand{\tightlist}{\setlength{\itemsep}{0pt}\setlength{\parskip}{0pt}}
\setlength{\emergencystretch}{3em}
\DeclareUnicodeCharacter{2212}{\ensuremath{-}}
\DeclareUnicodeCharacter{00D7}{\ensuremath{\times}}
\DeclareUnicodeCharacter{2265}{\ensuremath{\ge}}
\DeclareUnicodeCharacter{2264}{\ensuremath{\le}}
\DeclareUnicodeCharacter{2260}{\ensuremath{\ne}}
\DeclareUnicodeCharacter{2208}{\ensuremath{\in}}
\DeclareUnicodeCharacter{2209}{\ensuremath{\notin}}
\DeclareUnicodeCharacter{2227}{\ensuremath{\land}}
\DeclareUnicodeCharacter{2228}{\ensuremath{\lor}}
\DeclareUnicodeCharacter{222A}{\ensuremath{\cup}}
\DeclareUnicodeCharacter{2216}{\ensuremath{\setminus}}
\DeclareUnicodeCharacter{03C3}{\ensuremath{\sigma}}
\DeclareUnicodeCharacter{03C0}{\ensuremath{\pi}}
\DeclareUnicodeCharacter{03B5}{\ensuremath{\epsilon}}
\DeclareUnicodeCharacter{03B2}{\ensuremath{\beta}}
\DeclareUnicodeCharacter{2115}{\ensuremath{\mathbb{N}}}
\DeclareUnicodeCharacter{211D}{\ensuremath{\mathbb{R}}}
\DeclareUnicodeCharacter{2080}{\ensuremath{_0}}
\DeclareUnicodeCharacter{2081}{\ensuremath{_1}}
\DeclareUnicodeCharacter{2082}{\ensuremath{_2}}
\DeclareUnicodeCharacter{2083}{\ensuremath{_3}}
\DeclareUnicodeCharacter{2084}{\ensuremath{_4}}
\DeclareUnicodeCharacter{2085}{\ensuremath{_5}}
\DeclareUnicodeCharacter{2086}{\ensuremath{_6}}
\DeclareUnicodeCharacter{2087}{\ensuremath{_7}}
\DeclareUnicodeCharacter{2088}{\ensuremath{_8}}
\DeclareUnicodeCharacter{2089}{\ensuremath{_9}}
\DeclareUnicodeCharacter{2070}{\ensuremath{^0}}
\DeclareUnicodeCharacter{00B9}{\ensuremath{^1}}
\DeclareUnicodeCharacter{00B2}{\ensuremath{^2}}
\DeclareUnicodeCharacter{00B3}{\ensuremath{^3}}
\DeclareUnicodeCharacter{2074}{\ensuremath{^4}}
\DeclareUnicodeCharacter{2075}{\ensuremath{^5}}
\DeclareUnicodeCharacter{2076}{\ensuremath{^6}}
\DeclareUnicodeCharacter{2077}{\ensuremath{^7}}
\DeclareUnicodeCharacter{2078}{\ensuremath{^8}}
\DeclareUnicodeCharacter{2079}{\ensuremath{^9}}
\DeclareUnicodeCharacter{207A}{\ensuremath{^+}}
\DeclareUnicodeCharacter{207B}{\ensuremath{^-}}
\DeclareUnicodeCharacter{02B0}{\ensuremath{^h}}
\DeclareUnicodeCharacter{230A}{\ensuremath{\lfloor}}
\DeclareUnicodeCharacter{230B}{\ensuremath{\rfloor}}
\DeclareUnicodeCharacter{2192}{\ensuremath{\to}}
\DeclareUnicodeCharacter{2205}{\ensuremath{\emptyset}}
\begin{document}
\section*{A 3×4 ordered binary removal
obstruction}\label{a-34-ordered-binary-removal-obstruction}

October 11, 2026. \textbf{Candidate result with a written all-depth
proof.} The exact finite checks use two implementations. A separate
agent reviewed the argument; this is not independent mathematical peer
review. This new host and main theorem are now checked in Lean,
including the full quantifiers. See the source-bound audit and
verification scope in Section 8.

The tree, threshold entries, mode roles, and arbitrary-repair sampler
are inherited from OpenAI's September 25 manuscript. Erebos previously
proved the 4×4 obstruction and exact repair equality. Our construction
changes the anchor encoding, retaining the globally ordered modes from
the preceding \href{proofs/optimized-order.md}{4×4 quantitative
refinement}. The new ingredient is one anchor row with nonuniform extra
ones, described explicitly below.

\subsection*{Statement}\label{statement}

The fixed rectangular pattern is

\[
 H=\begin{pmatrix}1&0&0&1\\1&0&1&0\\0&1&1&1\end{pmatrix},
 \qquad P=\begin{pmatrix}1&0\\1&1\end{pmatrix}.
\]

For every integer h≥1, set m=2\^{}h, d=14h+2, n=dm. The square binary
host A\_h below satisfies

\[
 \operatorname{dist}_H(A_h)\ge d^{-2},\qquad
 N_H(A_h)\le 2mn^5,\qquad
 \frac{N_H(A_h)}{n^7}\le 2d^{-2}2^{-h}.
\]

Copies use three increasing distinct rows and four increasing distinct
columns and match all twelve entries. Distance allows arbitrary binary
changes in either direction anywhere in the host, normalized by n². No
equality for repair distance or exact all-depth copy count is claimed.

Consequently, for all c,C\textgreater0 and every ε₀\textgreater0, there
exists 0\textless ε\textless ε₀ such that hosts of arbitrarily large
order are ε-far from H-freeness and have fewer than c ε\^{}C n\^{}7
copies. In particular H admits no polynomial removal bound. This reduces
the pattern from the known 4×4 to 3×4 (16 to 12 entries); it does not
establish a 3×3 obstruction or minimality.

\subsection*{1. The host, including every binary
entry}\label{the-host-including-every-binary-entry}

Use a complete binary tree of height h, with nodes p=0,\ldots,2\^{}i−1
at depth i. For i\textless h put row blocks R\_i\^{}−(p),R\_i\^{}+(p)
and column blocks C\_i\^{}−(p),C\_i\^{}+(p), each of size m/2\^{}i. At
depth h identify the two signs: R\_h\^{}−(p)=R\_h\^{}+(p)=L\_R(p) and
C\_h\^{}−(p)=C\_h\^{}+(p)=L\_C(p), each of size one. Add dummy blocks
R\_\emph{,C\_} of size m. A class omits the node index and denotes the
union over its level; a parity subscript restricts node parity.

Rows follow recursive order: minus block, left subtree, right subtree,
plus block. Columns follow plus block, left subtree, right subtree,
minus block. A leaf block occurs once. Dummy rows are last, dummy
columns first among nonanchors. Positions within a block are twins.

Every nonanchor entry meeting a dummy is one. All variable entries are
zero except the following threshold entries:

\begin{longtable}[]{@{}lll@{}}
\toprule\noalign{}
Row / column & Value & Range \\
\midrule\noalign{}
\endhead
\bottomrule\noalign{}
\endlastfoot
R\_i\^{}+(p), C\_i\^{}+(q) & 1{[}p≤q{]} & 0≤i≤h \\
R\_i\^{}−(p), C\_i\^{}−(q) & 1{[}p\textless q{]} & 0≤i\textless h \\
R\_i\^{}+(p), C\_\{i−1\}\^{}+(q) & 1{[}⌊p/2⌋≤q{]} & 1≤i≤h \\
R\_i\^{}−(p), C\_\{i−1\}\^{}−(q) & 1{[}⌊p/2⌋\textless q{]} & 1≤i≤h \\
\end{longtable}

The last two rules both apply to a leaf row, with different column
signs. There is no conflicting minus diagonal rule at leaves.

There are six modes at each i=1,\ldots,h, with these disjoint roles:

\begin{longtable}[]{@{}lllll@{}}
\toprule\noalign{}
Mode & Row role 1 & Row role 2 & Column role 1 & Column role 2 \\
\midrule\noalign{}
\endhead
\bottomrule\noalign{}
\endlastfoot
V\_i\^{}+ & R\_i\^{}+ & R\_\{i−1\}\^{}+ & C\_* & C\_\{i−1\}\^{}+ \\
V\_i\^{}− & R\_\{i−1\}\^{}− & R\_i\^{}− & C\_* & C\_\{i−1\}\^{}− \\
W\_\{i,d\}\^{}+ & R\_\{i,d\}\^{}+ & R\_* & C\_\{i−1\}\^{}+ &
C\_\{i,d\}\^{}+ \\
W\_\{i,d\}\^{}− & R\_\{i,d\}\^{}− & R\_* & C\_\{i,d\}\^{}− &
C\_\{i−1\}\^{}− \\
\end{longtable}

For each mode t add \textbf{one row anchor block X\_t of size 2m} and
\textbf{two column anchor blocks Y\_\{t,0\},Y\_\{t,1\}, each of size m}.
Put anchors before all nonanchors. On both axes order modes as

\[
 (W^-_{1,0},W^-_{1,1},\ldots,W^-_{h,0},W^-_{h,1}),\quad
 (W^+_{1,0},W^+_{1,1},\ldots,W^+_{h,0},W^+_{h,1}),\quad
 (V^+_1,V^-_1,\ldots,V^+_h,V^-_h).
\]

Within each column mode, group 0 precedes group 1. Anchor row X\_t is
one on Y\_\{t,0\}, zero on every other anchor column. A nonanchor row in
row role j∈\{1,2\} of t has signature e\_j against
(Y\_\{t,0\},Y\_\{t,1\}); a row in neither role has signature 00.

It remains to specify each anchor row on nonanchor columns. Write
D=C\_*. Its one-set T\_t is exactly:

\begin{longtable}[]{@{}ll@{}}
\toprule\noalign{}
Anchor row mode & T\_t \\
\midrule\noalign{}
\endhead
\bottomrule\noalign{}
\endlastfoot
V\_i\^{}σ & C\_\{i−1\}\^{}σ \\
W\_\{i,d\}\^{}+ & D ∪ C\_\{i,d\}\^{}+ \\
W\_\{i,1\}\^{}− & D ∪ C\_\{i−1\}\^{}− \\
W\_\{i,0\}\^{}−, i\textless h & D ∪ C\_\{i−1\}\^{}− ∪ C\_\{i,1\}\^{}− \\
W\_\{h,0\}\^{}− & D ∪ C\_\{h−1\}\^{}− \\
\end{longtable}

In particular the odd-child addition is \textbf{absent at the common
leaves}. All unspecified entries are zero. The added dummy and odd-child
ones are outside the designated column roles. On those roles each X\_t
still has signature (0,1), as required by H's last two columns. Together
with the prescribed row signatures, every ordered designated P extends
to H.

Each axis has 12hm anchor positions and (2h+2)m nonanchor positions.
Thus the host is square with n=(14h+2)m.

\subsection*{2. Elementary support and order
facts}\label{elementary-support-and-order-facts}

At a fixed signed level, blocks occur in increasing node order. For a
fixed nonanchor row, entries along a column class are nondecreasing; for
a fixed nonanchor column, entries along a row class are nonincreasing.
Thus same-class rows or columns have nested neighborhoods.

Nonanchor rows and columns contain no unordered 3×3 crown J₃−I₃. Here is
the inherited argument. Same-class columns cannot occur in a crown. A
nondummy nonleaf row is supported on its own and preceding signed column
classes. A leaf row is supported on
L\_C,C\_\{h−1\}\^{}+,C\_\{h−1\}\^{}−. Hence three distinct classes with
pairwise common nondummy neighbors must be those three terminal classes.
On fixed columns from those classes, leaf traces are nested as the leaf
node increases, and nonleaf rows have zero on L\_C. They cannot provide
both required incomparable traces with a one on L\_C.

The one-set of any anchor column on nonanchor rows is a whole class or a
parity half-class. Any two such sets are disjoint or nested. No mode
order assumption is needed for this laminarity statement.

We repeatedly use the recursive-order comparisons, with π(p)=⌊p/2⌋:

\[
\begin{aligned}
 C_{i-1}^+(b)<C_i^\sigma(q)&\iff b\le\pi(q),&
 C_i^\sigma(q)<C_{i-1}^+(b)&\iff\pi(q)<b,\\
 C_i^\sigma(q)<C_{i-1}^-(b)&\iff\pi(q)\le b,&
 C_{i-1}^-(b)<C_i^\sigma(q)&\iff b<\pi(q),\\
 C_j^+(a)<C_j^-(b)&\iff a\le b,&
 C_j^-(a)<C_j^+(b)&\iff a<b,\\
 R_i^+(p)<R_{i-1}^+(a)&\iff\pi(p)\le a,&
 R_{i-1}^-(a)<R_i^-(p)&\iff a\le\pi(p).
\end{aligned}
\]

The two same-level cross-sign comparisons require j\textless h; at depth
h the signs are identified. Child signs in the parent/child comparisons
may be the common leaf sign. Reversing a strict block comparison gives
its complementary strict inequality. These facts follow directly from
the consecutive left-to-right subtree intervals.

\subsection*{3. Exactly one anchor row and one or two anchor
columns}\label{exactly-one-anchor-row-and-one-or-two-anchor-columns}

Consider a copy on rows r₁\textless r₂\textless r₃ and columns
c₁\textless c₂\textless c₃\textless c₄. Let a,b count the selected
anchors on each axis. Anchors are prefixes. H has three distinct rows
and four distinct columns, so no twin block can be reused.

\textbf{No anchor rows (a=0).} If b=0, columns 1,3,4 form a crown,
impossible. If b=1 or 2, column 1 places rows 1 and 2 in one class, but
their traces on nonanchor columns 3,4 are 01 and 10. If b≥3, anchor
columns 1 and 3 have traces 110 and 011, contradicting laminarity.

\textbf{At least two anchor rows with b\textgreater0.} Column 1 is one
on the first two anchor rows. It can only meet one mode's row block, so
those rows would be twins, impossible. Hence a≥2 requires b=0.

For the next cases let A\textless B be the modes of the first two anchor
rows. Their supports must have a common c₁ and exclusive columns
c₃∈T\_B∖T\_A, c₄∈T\_A∖T\_B. If c₁ is variable, c₄ cannot be dummy since
c₁\textless c₄. Inspection of the explicit support table leaves
precisely

\[
 A=W^-_{i,0},\quad i<h,\qquad
 B=W^-_{i+1,d}\ \text{or}\ V^-_{i+1}.
\]

Indeed the only other crossing pair with variable intersection is
W\_\{i,d\}\^{}+,V\_\{i+1\}\^{}+, whose first-only part is dummy.
Same-level minus supports are nested, disjoint-level supports have no
variable intersection, and vertical supports are single distinct
classes. For the displayed pairs, c₁ is in odd C\_i\^{}−, c₄ in
C\_\{i−1\}\^{}−, and c₃ in even C\_i\^{}−, or possibly odd
C\_\{i+1\}\^{}− when B=W\^{}-\_\{i+1,0\} with i+1\textless h.

\textbf{Two anchor rows (a=2,b=0).} Column c₁ is variable since its
bottom entry is zero. In the displayed pair classification, a nondummy
row cannot be one on both C\_\{i−1\}\^{}− and C\_\{i+1\}\^{}−. Therefore
c₃ is in even C\_i\^{}− and r₃ must be in R\_i\^{}−. Its entire
nonanchor one-support lies in D∪C\_\{i−1\}\^{}−∪C\_i\^{}−, covered by
T\_A∪T\_B. This contradicts c₂'s trace 001.

\textbf{Three anchor rows (a=3,b=0).} If c₁ is variable, the same pair
classification applies. No anchor row is one both on c₄∈C\_\{i−1\}\^{}−
and c₃ in even C\_i\^{}− or odd C\_\{i+1\}\^{}−: the sole support
spanning the adjacent minus classes adds only odd children. If c₁ is
dummy, its trace 110 forces the first two modes to be W and the last to
be V. The V support is one full class K, so c₂,c₃,c₄ all lie in K. A W
support restricted to K is empty, all of K, or one parity half. The
exclusive c₃,c₄ force opposite halves, whose union covers K; there is no
c₂ with trace 001.

It follows that a=1. If b=4, c₁ and c₄ must both be the sole group-0
anchor class on which r₁ is one; these are twins, impossible. If b=3,
column c₃ places r₂,r₃ in one row class, whereas the nonanchor c₄ has
trace 01 on them, contradicting nonincreasing same-class row entries.

\textbf{No anchor columns (a=1,b=0).} Column c₁ is not dummy, since its
bottom entry is zero; c₃ is not dummy either, since
c₁\textless c₂\textless c₃. Both lower rows are nondummy and lie in
different classes, since their column traces include both 10 and 01.
Their common one at c₃ forces adjacent signed row classes, possibly
including the terminal leaf class.

For a V mode, a W\^{}+ mode, or an unextended W\^{}− mode, c₁,c₄ lie in
one column class. Their order contradicts r₂·c₁=1,r₂·c₄=0. The sole
remaining mode is W\^{}-\emph{\{i,0\}, i\textless h. Put
P₀=C}\{i−1\}\^{}− and Q₀=odd C\_i\^{}−; c₁,c₄ must lie in different
classes.

If c₁∈P₀ and c₄∈Q₀, the lower row classes are either
(R\_\{i−1\}\^{}−,R\_i\^{}−), forcing their common c₃ into P₀ where its
top entry is one, or (R\_i\^{}−,R\_\{i+1\}\^{}−). In the latter case,
write c₁=P₀(b), c₄=Q₀(q), r₂=R\_i\^{}−(p). Column order gives
b\textless π(q), while r₂·c₁=1 gives π(p)\textless b. Hence
p\textless q, contradicting r₂·c₄=0.

If c₁∈Q₀ and c₄∈P₀, the lower row classes are either
(R\_i\^{}−,R\_\{i−1\}\^{}−), again forcing c₃ into P₀, or
(R\_\{i+1\}\^{}−,R\_i\^{}−). In the latter case let r₂ have node p and
r₃ node z. Since c₂ has top zero and bottom one, it must be an even
C\_i\^{}− column with node u: dummy and P₀ are in the top support. Row
order gives π(p)\textless z; the bottom entry gives z\textless u. Thus
r₂·c₂=1, contradicting its required zero. If i+1=h, the same minus
incidences and row-order comparison hold with the common leaf class.
This finishes the b=0 exclusion.

We have proved (a,b)∈\{(1,1),(1,2)\}.

\subsection*{4. Every copy uses its own terminal minus
parent}\label{every-copy-uses-its-own-terminal-minus-parent}

Let t be r₁'s mode. Column c₁ is Y\_\{t,0\}, so r₂ belongs to row role 1
of t. The remaining conditions include

\[
 c_3<c_4,\quad A(r_2,c_3)=1,\quad A(r_2,c_4)=0,
 \quad c_3\notin T_t,\quad c_4\in T_t.
\]

Both c₃,c₄ are nonanchors. Column c₄ is not dummy, since c₃\textless c₄.

\textbf{Vertical t.} For V\_i\^{}+, r₂ is R\_i\^{}+(p) and
c₄=C\_\{i−1\}\^{}+(q) with π(p)\textgreater q. A nondummy c₃ outside
T\_t with a one on r₂ would be in C\_i\^{}+ (or, at a leaf,
C\_\{h−1\}\^{}−). Its order before c₄ contradicts the threshold
inequalities. Explicitly a C\_i\^{}+ node u requires p≤u and
π(u)\textless q; a terminal minus-parent node u requires
π(p)\textless u\textless q. For V\_i\^{}−, r₂ is R\_\{i−1\}\^{}−(p) and
c₄=C\_\{i−1\}\^{}−(q) with p≥q. A nondummy c₃ outside T\_t could only be
C\_\{i−2\}\^{}−(u) when i≥2, requiring π(p)\textless u\textless π(q),
also impossible. At i=1 that preceding class is absent. Thus c₃ must be
dummy.

When b=1, the nonanchor c₂\textless c₃ rules that out. When b=2, a later
nondummy r₃ with a one on c₄ is impossible by the same threshold and
row-order inequalities: in the plus case its row class is
R\_\{i−1\}\^{}+ or R\_i\^{}+, and in the minus case R\_\{i−1\}\^{}− or
R\_i\^{}−. Hence r₃ is dummy. Its one at anchor c₂ would require a W
group-1 column. Every W anchor precedes the V column c₁, contradicting
c₁\textless c₂.

\textbf{Positive horizontal t=W\_\{i,d\}\^{}+.} Column c₄ is in
C\_\{i,d\}\^{}+, as is the class of r₂. Write their nodes p,q. The zero
requires p\textgreater q. A c₃ in the same child class is excluded by
monotonicity and c₃\textless c₄. A plus parent c₃=C\_\{i−1\}\^{}+(u)
requires π(p)≤u≤π(q). Equal floors would make p,q unequal siblings,
contradicting their common parity d.~At the terminal level a
minus-parent c₃ instead requires π(p)\textless u\textless π(q),
impossible. The dummy is in T\_t and cannot be c₃. No copy occurs in
this mode.

\textbf{Negative horizontal t=W\_\{i,d\}\^{}−.} The row r₂ is
R\_\{i,d\}\^{}−(p), or L\_R(p) if i=h. If c₄ is an added odd child, then
i\textless h,d=0. Its node q satisfies p≥q. Column c₃ outside T\_t with
a one on r₂ can only be an own-level node u\textgreater p;
c₃\textless c₄ would give u\textless q, impossible.

Thus c₄=C\_\{i−1\}\^{}−(b), with π(p)≥b. For i\textless h, a c₃ outside
T\_t with a one on r₂ must be C\_i\^{}−(q), where q\textgreater p and
π(q)≤b. These inequalities force π(q)=b=π(p), p even and q=p+1 odd. The
row role then has d=0, but this odd c₃ was added to T\_t, a
contradiction. This is the precise use of the new odd-child ones.

Therefore i=h. Here c₃ can only be a leaf column q with p≤q, or a
plus-parent column C\_\{h−1\}\^{}+(a) with π(p)≤a. In the first case its
order before c₄ gives π(q)≤b; in the second it gives a≤b. Together with
π(p)≥b, both cases imply b=π(p).

Every copy consequently sends its fixed zero H(2,4) to

\[
 T_h=\{(L_R(p),c):c\in C_{h-1}^-(\lfloor p/2\rfloor)\}.
\]

There are exactly 2m cells in T\_h, because each terminal parent block
has two positions. Fixing this cell leaves at most n⁵ choices for the
other two rows and three columns. Thus N\_H(A\_h)≤2mn⁵.

\subsection*{5. Arbitrary repairs still cost at least
m²}\label{arbitrary-repairs-still-cost-at-least-muxb2}

Let a binary B differ from A\_h on M\textless m² cells. Sample a uniform
leaf z and its ancestor at each level. Independently choose one uniform
position from each signed row and column block on this path, identifying
the two sign selections at the leaf. Choose independent uniform
representatives of each anchor block and each dummy block.

A variable or dummy position has marginal probability 1/m. A column
anchor position has probability 1/m, and a row anchor position 1/(2m).
Selections with an anchor or dummy endpoint are independent across axes;
root selections are also independent. Therefore any protected cell
(anchor endpoint, dummy endpoint, or root-to-root) is selected with
probability at most 1/m². The union bound supplies a selection with no
edited protected cell. This probability assertion is not made for
arbitrary variable-variable cells or for shared leaf diagonals.

Suppose B is H-free. The unchanged designated anchors extend every
ordered P body on the selected role representatives to H. Therefore all
those bodies must avoid P. The root-plus diagonal is one and the
root-minus diagonal zero. If the selected parent diagonals have these
values, the vertical bodies are

\[
 \begin{pmatrix}1&b_+\\1&1\end{pmatrix},\qquad
 \begin{pmatrix}1&0\\1&b_-\end{pmatrix}.
\]

Their P-freeness forces b\_+=1,b\_−=0. The horizontal modes matching the
child parity next have bodies

\[
 \begin{pmatrix}1&c_+\\1&1\end{pmatrix},\qquad
 \begin{pmatrix}c_-&0\\1&1\end{pmatrix},
\]

forcing child diagonal values c\_+=1,c\_−=0. The recursive axis orders
put all these role bodies in increasing order. Induction reaches the
common leaf diagonal, which cannot simultaneously be one and zero. Thus
B is not H-free. All M edits were arbitrary, including anchor edits and
both directions of binary change. This proves the m² lower bound.

\subsection*{6. Quantifiers and
comparison}\label{quantifiers-and-comparison}

With ε\_h=(14h+2)\^{}−2, the density-to-power ratio is at most

\[
 \frac{N_H(A_h)}{n_h^7\epsilon_h^C}
 \le 2(14h+2)^{2C-2}2^{-h}\longrightarrow0.
\]

Given c,C,ε₀\textgreater0, choose h with ε\_h\textless ε₀ and this
ratio\textless c.~Replace every physical row and column by q identical
consecutive twins, for any positive integer q. The same sampler gives
repair lower bound q²m². Since H has distinct rows and columns, a copy
cannot reuse a projected row or column, so its count scales exactly by
q⁷. Both normalized bounds are unchanged. Choosing q arbitrarily large
gives the required host sizes, including a prescribed threshold
depending on ε\_h.

The improvement is rectangular pattern size: 3×4 rather than the known
4×4. Copy densities have different normalizations n⁷ and n⁸, so the raw
count is not a multiplicative improvement to the old fixed-H bound. The
earlier 4×4 exact-distance and optimized-count results remain separate.
Neither the exact repair equality nor its existing Lean certificate has
been transferred to this host.

\subsection*{7. Reproduction and verification
scope}\label{reproduction-and-verification-scope}

\texttt{experiments/asymmetric\_candidate.py} enumerates all ordered
quotient copies with exact twin multiplicities through h=6.
\texttt{verification/check\_three\_by\_four.py} separately constructs
the host and uses weighted subsequence dynamic programming. It checks
the same depths, a fully expanded h=1 matrix, designated anchor
signatures, and adverse variants omitting each of the two anchor-row
changes. The latter variants have nonlocalized copies. See the archived
exact report and separate agent review. These finite computations
support this written argument; they do not supply its universal
quantifiers or independent peer review.

\subsection*{8. Lean formalization}\label{lean-formalization}

The main theorem is checked in Lean 4.34.1 with the pinned Mathlib
dependency.

\texttt{no\_polynomial\_removal\_bound} in the
\texttt{OrderedRemoval.Rectangular} namespace states the full quantifier
order: for every c,C,ε₀\textgreater0, choose a fixed
0\textless ε\textless ε₀; for every host threshold N, there is an
n≥max(1,N) square binary matrix with at least εn² edits in every H-free
repair and fewer than c ε\^{}C n⁷ copies.

The proof constructs the literal host and both orders, excludes every
unwanted anchor configuration, localizes the fixed zero, counts copies,
and proves the arbitrary-repair bound before normalizing. The
uniform-blowup lemmas preserve the normalized bounds. They prove the
sufficient q⁷ copy upper bound; the exact blowup-count equality in
Section 6 is not needed by the formal capstone. No exact all-depth
enumeration or repair equality is asserted for this host.

The formal sampler is an alternative to Section 5's independent sampler:
it retains the inherited three seeds and selects representatives from
both m-sized row subgroups of each 2m anchor block. Copy completion uses
the second representative. Its anchor-row marginal is 1/m, and every
protected cell still has probability at most 1/m². This distinction
changes no bound.

The \href{verification/reference/rectangular-formal-audit.json}{audit}
binds all project source files. It checks transitive axiom dependencies
and replays each project module into its imported Lean environment. This
does not freshly replay all Mathlib or constitute independent
mathematical peer review. Only \texttt{propext}, \texttt{Quot.sound},
and \texttt{Classical.choice} occur as axioms. The separate-agent
semantic review is archived alongside the finite checks.

\end{document}
